\documentclass[10pt,a4paper]{article}
\usepackage{amsmath,amssymb,amsthm,graphicx}
\title{Conditional hardness of algebraic-series truncation via integer factorization and ultimate factorial complexity}
\author{Pranjal Dutta, Gorav Jindal, Anurag Pandey, Amit Sinhababu}
\date{doi:10.4230/LIPIcs.CCC.2021.25; CC BY4.0}
\begin{document}
\maketitle
\noindent\textit{Editorial note.} Original authors: Pranjal Dutta, Gorav Jindal, Anurag Pandey, Amit Sinhababu. Official CCC2021 article25, DOI10.4230/LIPIcs.CCC.2021.25.\par\medskip
\noindent\textit{Editorial note.} License: https://creativecommons.org/licenses/by/4.0/\par\medskip
\noindent\textit{Editorial note.} Faithful editable TeX transcription of formal Theorems32/35, complete author Sections6.1/6.2 proofs, full integral-series Theorem77 proof and AppendixI Algorithm1. Original constant-free/nonuniform/fixed-k/all-i/large-d conditions, author expressions and complete published bibliography retained without mathematical correction or new proofs. Only official CC BY4.0 version used; sketch-only permanent consequence excluded.\par\medskip
\noindent\textit{Allowed standard prerequisites (editorial).}\begin{itemize}
\item Binomial series and integer divisibility. Exact author binomial-series formula, Fact74 and Legendre Theorem76 retained.
\end{itemize}\medskip
\section*{Author integer-complexity and integrality context}
\paragraph{Definition 11 (Straight Line Program).} An SLP (straight line program) $P$ (for computing an integer) of length (or size) $n$ is a sequence of integers $a_0,\ldots,a_n$ with $a_0=1$ and $a_k=a_i\circ a_j$ with $i,j<k$ for $\circ\in\{+,-,\times\}$. We say that the SLP $P$ computes the integer $a_n$. For an integer $N$, we define the straight line complexity $\tau(N)$ of $N$ to be the length of the smallest SLP computing $N$.
\paragraph{Constant-free circuit complexity.} In the same spirit, one can define constant-free circuit complexity of polynomials where the given constants belong to the set $\{-1,0,1\}$.\footnote{To use $2^n$ in the circuit, one has to build up a circuit for $2^n$, of size $\log n$, from 1; whereas in the usual sense of circuit size, constants are free. Thus, $f_d:=2^{2^d}x^d$ has $O(\log d)$-size circuit but requires $\Omega(d)$-size constant-free circuit.} We denote, $\tau(f)$ as the size of the minimal constant-free circuit computing $f$. Trivially, $L(f)\le\tau(f)$.
\paragraph{Definition 28 (Easy sequence).} A sequence $(a_n)_n$ of integers is said to be ``easy to compute'' if there exists a polynomial $p$ such that straight line complexity of $a_n$, i.e. $\tau(a_n)\le p(\log n)$, for $n\ge1$. If a sequence is not easy to compute, it is said to be hard. In fact, for most numbers $N$, one can show that $\tau(N)\ge\log N/\log\log N$ (``close'' to the trivial upper bound) [14,35]. It is believed that $(d!)$ is hard to compute. In fact, its hardness is deeply connected to the infamous integer factoring problem. [42] showed that $d!$ being easy to compute will imply factoring is easy in the non-uniform setting.\footnote{However, this result does not imply that natural numbers can be factored in polynomial time in the Turing-Machine model, as the numbers used can be $\mathrm{poly}(n)$-bits.}
\paragraph{Definition 33 (Ultimately easy).} A sequence of integers $(a_n)$ is ultimately easy if there exists another sequence $(b_n)$ such that $\tau(a_nb_n)\le\mathrm{poly}(\log n)$ for all large enough $n$.
\paragraph{Definition 34 (Ultimate complexity).} Define the ultimate complexity of an integer $n$ as the minimum $\tau$-complexity of its multiple, i.e. $\tau_1(n)=\min_{b\in\mathbb Z\setminus\{0\}}\tau(b\cdot n)$. It is clear that Definition 33 can be stated wrt $\tau_1$. We remark that $\tau_1(n_1n_2)\le\tau_1(n_1)+\tau_1(n_2)+1$, for any $n_1,n_2\in\mathbb Z$.
\paragraph{Binomial series.} For rational $n$, $(x+a)^n=\sum_{k\ge0}\binom nkx^ka^{n-k}$, where $\binom nk=n(n-1)\cdots(n-k+1)/k!$. A formal power series is $A=\sum_{i\ge0}A_ix^i$; its degree-$d$ truncation is $\mathrm{trunc}(A,d)=\sum_{0\le i\le d}A_ix^i$.
\paragraph{Fact 74 (Folklore).} Product of any $k$ consecutive positive integers is divisible by $k!$.
\paragraph{Definition 75 ($p$-adic valuation).} Let $p$ be a prime and $n\in\mathbb Z$. We denote $p$-adic valuation of $n$ as $\nu_p(n)$ to be the highest exponent such that $p^{\nu_p(n)}\mid n$. Formally, $\nu_p:\mathbb Z\longrightarrow\mathbb N$ defined by $\nu_p(n)=\max\{v\in\mathbb N:p^v\mid n\}$. Note that, by definition, $\nu_p(\mathrm{rad}(n))=1$ if $p\mid n$, and 0 otherwise.
\paragraph{Theorem 76 (Legendre's formula).} For a prime $p$ and $n\in\mathbb N$, $\nu_p(n)=\sum_{j=1}^\infty\lfloor n/p^j\rfloor$.
\paragraph{Theorem 77 (Restatement of Theorem 31, Integral power series).} Let $k\in\mathbb N$. Define $f_k(x):=(1+k^2x)^{1/k}$. Then, $f_k(x)\in\mathbb Z[[x]]$.
\begin{proof}By binomial expansion, $f_k\in\mathbb Q[[x]]$. Let $f_k(x)=\sum_{d\ge0}a_dx^d$. We'll prove by strong induction that indeed the coefficients are integers. Obviously $a_0=1$, and assume that for $m\in\mathbb N$ we have proved that $a_\ell\in\mathbb Z$ for $0\le\ell<m$. The coefficient at $x^m$ in $(\sum_{d=0}^\infty a_dx^d)^k=(1+\sum_{d=1}^\infty a_dx^d)^k$ is equal to $k\cdot a_m$ plus a bunch of terms that we know are integer by the induction hypothesis; hence $k\cdot a_m=b\in\mathbb Z$. But by the binomial series formula we have
\[a_m=k^{2m}\binom{1/k}m=\frac{k^{2m}\prod_{j=0}^{m-1}(1/k-j)}{m!}=\frac{k^m\prod_{j=0}^{m-1}(1-kj)}{m!}.\]
It suffices to prove that $k\mid b$. If we can show that $\nu_p(b)\ge\nu_p(k)$ for every prime $p$ dividing $k$, this would certainly imply that $k\mid b$. So, fix a prime $p\mid k$. Note that
\[b=k\cdot a_m=X/m!,\qquad X:=k^{m+1}\prod_{j=0}^{m-1}(1-kj).\]
As, $p\mid k$, we must have $\prod_{j=0}^{m-1}(1-kj)\equiv1\bmod p$. Thus, $\nu_p(X)=\nu_p(k^{m+1})=(m+1)\nu_p(k)$. And by Theorem 76, $\nu_p(m!)=\sum_{j=1}^\infty\lfloor m/p^j\rfloor<\sum_{j=1}^\infty m/p^j=m/p-1\le m$. Thus,
\[\nu_p(b)=\nu_p(X)-\nu_p(m!)\ge(m+1)\nu_p(k)-m\ge\nu_p(k),\]
as we wanted. Putting it together gives $a_m\in\mathbb Z$ proving the inductive step. Hence, the conclusion follows.\end{proof}
It suffices to prove the integrality of $(1+k^2x)^{1/k}$ for the candidate functions $(1+k^2x)^{i/k}$ with $i<k$, as in the authors' Appendix G. In Section6 the authors work throughout with $(1+k^2x)^{i/k}$, $i,k\in\mathbb N$ and $i<k$; Theorem31 states their integrality.

\section*{Formal author conditional algebraic-truncation results}
\paragraph{Theorem 32.} Let $k\in\mathbb N$. If $\tau(\mathrm{trunc}((1+k^2x)^{i/k},d))=O(\log^cd)$ (for some constant $c$) for all $i\in[k-1]$ then integer factorization (in the non-uniform setting) can be performed in polynomial time.
\paragraph{Theorem 35.} Fix $k\in\mathbb N$. Suppose, for each $i\in[k-1]$, there exists some constant $c$ such that $\tau(\mathrm{trunc}((1+k^2x)^{i/k},d))=O(\log^cd)$, for large enough $d$. Then, $(n!)_{n\in\mathbb N}$ is ultimately easy.

\section*{Complete author Section 6.1 proof}
Here, we show that if the truncation of each $(1+k^2x)^{i/k}$, for $i\in[k-1]$, has small constant-free circuit, then one can factor $n$ in $\mathrm{poly}(\log n)$ time, in the non-uniform setting. This would readily imply the first part of Theorem 8.
\paragraph{Proof of Theorem 32.} Let, $(1+k^2x)^{i/k}=\sum_{d\ge0}C_{d,i}x^d\in\mathbb Z[[x]]$, where the coefficient $C_{d,i}$ of $x^d$ is equal to $\pm k^d(-i)(k-i)(2k-i)\cdots((d-1)k-i)/d!$. We see that the product of all $C_{d,i}$ is equal to:
\[\prod_{i\in[k-1]}C_{d,i}=\pm\frac{k^{(k-1)d}(k-1)!(dk)!}{(d!)^kk^d(kd-1)(kd-2)\cdots(kd-(k-1))}.\]
The assumption $\tau(\mathrm{trunc}((1+k^2x)^{i/k},d))=O(\log^cd)$ implies that $\tau(C_{d,i})=O(\log^cd)$ (just by subtracting two consecutive truncations and substituting $x=1$). This further implies that $\tau(\prod_{i\in[k-1]}C_{d,i})=O(\log^cd)$. Let us define, for any $d\ge1$,
\[N(d,k):=\frac{k^{(k-2)d}(dk)!}{(d!)^k}.\]
We first argue that $N(d,k)\in\mathbb N$. This follows from the fact that $N(d,k)=\prod_{i\in[k-1]}C_{d,i}\cdot(kd-1)\cdots(kd-(k-1))/(k-1)!$, and $(k-1)!$ must divide $(kd-1)\cdots(kd-(k-1))$, by Fact 74. Further, since $k$ is constant, it implies that $\tau(N(d,k))=O(\log^cd)$ (because the extra term has trivial $O(\log d)$-complexity).
Now, we describe how to find a non-trivial factor of a given integer $n$. We assume that all the primes dividing $n$ are larger than $k$; otherwise we can remove all the prime factors smaller than $k+1$ (since $k$ is a constant).
The idea is to first find a positive integer $t$ such that all the primes dividing $n$ are in the interval $[t+1,tk]$, by using an iterative algorithm; if such a $t$ does not exist we would have already found a non-trivial factor of $n$ (by the algorithm). As an invariant, we maintain an integer $m$ such that all the prime divisors of $n$ are greater than $m$. We start with $m=k$ and compute $\gcd(N(m,k),n)$ at each iteration. Since all the primes dividing $n$ are greater than $m$ (by assumption), we get that $\gcd(N(m,k),n)=\gcd((mk)!,n)$. If the $\gcd((mk)!,n)\ne1,n$, we must have already found a non-trivial factor of $n$ and we are done. Otherwise, we can have two cases: either (i) $\gcd((mk)!,n)=1$, or (ii) $\gcd((mk)!,n)=n$.
If $\gcd((mk)!,n)=1$ then we set $m\leftarrow mk$ and continue (because in this case all the primes dividing $n$ must be greater than $mk$). Otherwise we have $\gcd((mk)!,n)=n$, and hence, all the primes dividing $n$ are in the interval $[m+1,mk]$ and we stop with $t\leftarrow m$. We know that $t\le\lceil n/k\rceil$ and this uses at most $\log_kn=\log n$ iterations. So, this step has given us an integer $t$ such that all the primes dividing $n$ are in the interval $[t+1,tk]$, and the time taken is $\mathrm{poly}(\log n)$, due to only $\log n$ many iterations and each step takes $\mathrm{poly}(\log n)$-time due to the fact that $\tau(N(d,k))=O(\log^cd)$ implies gcd computation can be done in $\mathrm{poly}(\log n)$ (by euclidean algorithm).
Once, we know that all the primes are in an interval of the form $[t+1,tk]$, we now try to reduce the length of it to $k$ so that, we can simply brute force to get a factor of $n$, otherwise of course our algorithm would already find a factor. The length reduction part is similar to binary search algorithm that we describe below.
To find a positive integer $s$ such that all the primes dividing $n$ are in the interval $[sk+1,(s+1)k]$ (Or we find a non-trivial factor of $n$), again we use an iterative algorithm. As an invariant, we maintain two positive integers $a,b$ such that all the prime divisors of $n$ are in the interval $[ak+1,bk]$. We start with $a=1,b=t$. Our invariant is trivially true at the start.
At each iteration, we set $c=\lceil(a+b)/2\rceil$ and compute $\gcd(N(c,k),n)$. Since $c\le t$ and all the prime divisors of $n$ are larger than $t$, we get that $\gcd(N(c,k),n)=\gcd((ck)!,n)$. Again, we argue in the same way as before. If the gcd is $\ne1,n$, we have already found a non-trivial factor of $n$ and we are done. Otherwise, we have two cases: either (i) $\gcd((ck)!,n)=1$, or (ii) $\gcd((ck)!,n)=n$.
If $\gcd((ck)!,n)=1$ then it is clear that all the primes dividing $n$ are in the interval $[ck+1,bk]$ and hence we set $a\leftarrow c,b\leftarrow b$. If $\gcd((ck)!,n)=n$ then they all the primes dividing $n$ are in the interval $[ak+1,ck]$ and hence we set $a\leftarrow a,b\leftarrow c$. This will terminate when $b-a\le1$. Hence we find the desired positive integer $s$. This uses at most $\log t=\log n$ iterations.
Now we just need to search for the prime divisors of $n$ in the interval $[sk+1,(s+1)k]$ (an interval of constant length). Now, we brute force to finally find a non-trivial factor of $n$. Similarly, this step also takes $\mathrm{poly}(\log n)$ as each gcd computation takes $\mathrm{poly}(\log n)$ time. So, we have successfully found a non-trivial factor of $n$ by the end of this process, repeating this, we can get all the factors in $\mathrm{poly}(\log n)$-time and we are done.\qed
We also refer to Algorithm 1 in Appendix I, reproduced below.
\section*{Complete author Section 6.2 proof}
In this section, we show that easiness of truncation of $(1+k^2x)^{i/k}$ shows that a multiple of $n!$ must be easy. Note that, this may not imply that $n!$ is easy, however, from complexity-theoretic point-of-view, it is believed to be hard because of non-trivial implications. Shub \& Smale [44] proved: If $n!$ is ultimately hard to compute, then $\mathrm P\ne\mathrm{NP}$ over the field of complex numbers. Here, the computation is over Blum-Shub-Smale (BSS) model and can use complex numbers in the algorithm. In fact, a stronger version (known as $\tau$-conjecture) connects $z(f)$, distinct integer roots of $f$ with $\tau(f)$. Recently, [16] showed that a similar conjecture, in the SOS-model, would in fact imply explicit constructions of rigid matrices \& $\mathrm{VP}\ne\mathrm{VNP}$. For similar related works, we refer to [25,27].
\paragraph{Proof of Theorem 35.} Let, $(1+k^2x)^{i/k}=\sum_{d\ge0}C_{d,i}x^d\in\mathbb Z[[x]]$. From the hypothesis, it follows that there exists $c$ such that $\tau(C_{d,i})\le\log^cd$, for each $i\in[k-1]$ (subtract two consecutive terms and substitute $x=1$). Further, from the proof in Section 6.1 (and following the same notation), we know that
\[\prod_{i\in[k-1]}C_{d,i}=\pm\frac{k^{(k-2)d}(k-1)!(dk)!}{(d!)^k(kd-1)(kd-2)\cdots(kd-(k-1))}.\]
Let us define, $a(d,k):=(dk)!/(d!)^k$. Note that, $a(d,k)\in\mathbb N$ (it is the multinomial coefficient $\binom{dk}{d,\ldots,d}$). Further, $k^{(k-2)d}a(d,k)=\prod_{i\in[k-1]}C_{d,i}\cdot(kd-1)\cdots(kd-(k-1))/(k-1)!$, and $(k-1)!$ must divide $(kd-1)\cdots(kd-(k-1))$, by Fact 74. As $k$ is constant, each $kd-i$ can be computed in $O(\log d)$-time trivially. Further, $\tau(\prod_{i\in[k-1]}C_{d,i})\le O(\log^cd)$. As $\tau$ is additive over multiplication, it follows that
\[\tau(k^{(k-2)d}a(d,k))\le O(\log^cd)\implies\tau_1(a(d,k))\le O(\log^cd).\]
Now we recurse by noticing the following trivial identity that $n!=n!/(\lfloor n/k\rfloor)!)^k\cdot((\lfloor n/k\rfloor)!)^k$. We know by the above relation on $a(d,k)$ (and replacing $d:=\lfloor n/k\rfloor$ for some integer $n$) that
\[\tau_1\left(\frac{(k\lfloor n/k\rfloor)!}{(\lfloor n/k\rfloor)!^k}\right)\le O(\log^cn).\]
Further, any integer $n$ can be written as $n=k\lfloor n/k\rfloor+j$ for some $j\le k-1$. Note that $k\lfloor n/k\rfloor+j$ has complexity at most $\log n$ for each $j\in[k-1]$. So, multiplying $k\lfloor n/k\rfloor+j$ for $j\in[k-1]$, it is straightforward to deduce that
\[\tau_1\left(\frac{n!}{(\lfloor n/k\rfloor!)^k}\right)\le O(\log^cn).\tag{2}\]
As, $n!=n!/(\lfloor n/k\rfloor!)^k\cdot((\lfloor n/k\rfloor)!)^k$, and $\tau_1((\lfloor n/k\rfloor!)^k)\le\tau_1(\lfloor n/k\rfloor!)+O(1)$; use Equation (2):
\begin{align*}\tau_1(n!)&\le\tau_1(\lfloor n/k\rfloor!)+O(\log^cn)+O(1)\\&\le\tau_1(\lfloor n/k^2\rfloor!)+O(\log^cn)+O(\log^cn)+O(1)\\&\ \vdots\\&\le\log_kn\cdot O(\log^cn)=O(\log^{c+1}n).\end{align*}
Therefore, $(n!)$ is ultimately easy to compute, as we wanted.\qed

\section*{Author Appendix I, Algorithm 1}
Integer factorization assuming the truncations of $(1+k^2x)^{i/k}$ being easy for each $i$. Input: A composite positive integer $n$. Output: A non-trivial factor of $n$.
\begin{enumerate}
\item Define $N(d,k):=k^{(k-2)d}(dk)!/(d!)^k$.
\item $m\leftarrow k$.
\item \textbf{while true do}
\item Compute $\gcd(N(m,k),n)$.
\item \textbf{if} $\gcd(N(m,k),n)=1$ \textbf{then}
\item $m\leftarrow mk$.
\item \textbf{else if} $\gcd(N(m,k),n)=n$ \textbf{then}
\item $t\leftarrow m$. This $m$ is the desired $t$.
\item \textbf{break}
\item \textbf{else}
\item \textbf{return} $\gcd(N(m,k),n)$. Here $\gcd(N(m,k),n)$ is a non-trivial factor of $n$.
\item \textbf{end if}
\item \textbf{end while}. At this step, all the primes dividing $n$ are in the interval $[t+1,tk]$.
\item $a\leftarrow1$.
\item $b\leftarrow t$.
\item \textbf{while true do}
\item \textbf{if} $b-a\le1$ \textbf{then}
\item $s\leftarrow a$. This $a$ is the desired $s$.
\item \textbf{break}
\item \textbf{end if}
\item $c\leftarrow\lceil(a+b)/2\rceil$.
\item Compute $\gcd(N(c,k),n)$.
\item \textbf{if} $\gcd(N(c,k),n)=1$ \textbf{then}
\item $a\leftarrow c$.
\item \textbf{else if} $\gcd(N(c,k),n)=n$ \textbf{then}
\item $b\leftarrow c$.
\item \textbf{else}
\item \textbf{return} $\gcd(N(c,k),n)$. Here $\gcd(N(c,k),n)$ is a non-trivial factor of $n$.
\item \textbf{end if}
\item \textbf{end while}. At this step, all the primes dividing $n$ are in the interval $[sk+1,(s+1)k]$.
\item \textbf{for} $i=sk+1$ \textbf{to} $(s+1)k$ \textbf{do}
\item \textbf{if} $i$ divides $n$ \textbf{then}
\item \textbf{return} $i$. Here $i$ is a non-trivial factor of $n$.
\item \textbf{end if}
\item \textbf{end for}
\end{enumerate}

\begin{thebibliography}{99}
\bibitem{ref1} Alexander Alder. Grenzrang und Grenzkomplexität aus algebraischer und topologischer Sicht. PhD thesis, Zentralstelle der Studentenschaft, 1984.
\bibitem{ref2} Eric Allender, Peter Bürgisser, Johan Kjeldgaard-Pedersen, and Peter Miltersen. On the Complexity of Numerical Analysis. SIAM Journal on Computing, 38, January 2006. Preliminary version in the 21st Annual IEEE Conference on Computational Complexity (CCC’06). doi: 10.1109/CCC.2006.30.
\bibitem{ref3} Robert Andrews. Algebraic Hardness Versus Randomness in Low Characteristic. In 35th Computational Complexity Conference (CCC 2020), volume 169 of Leibniz International Proceedings in Informatics (LIPIcs), pages 37:1–37:32. Schloss Dagstuhl–Leibniz-Zentrum für Informatik, 2020. doi:10.4230/LIPIcs.CCC.2020.37.
\bibitem{ref4} Jean Berstel and Christophe Reutenauer. Rational Series and Their Languages. Springer- Verlag, Berlin, Heidelberg, 1988. URL: https://dl.acm.org/doi/book/10.5555/52107.
\bibitem{ref5} Markus Bläser and Gorav Jindal. On the Complexity of Symmetric Polynomials. In 10th Innovations in Theoretical Computer Science Conference (ITCS’19), volume 124 of LIPIcs, pages 47:1–47:14. Schloss Dagstuhl - Leibniz-Zentrum für Informatik, 2019. doi:10.4230/ LIPIcs.ITCS.2019.47.
\bibitem{ref6} Alin Bostan, Gilles Christol, and Philippe Dumas. Fast computation of the Nth term of an algebraic series over a finite prime field. In Proceedings of the ACM on International Symposium on Symbolic and Algebraic Computation (ISSAC’16), pages 119–126, 2016. doi: 10.1145/2930889.2930904.
\bibitem{ref7} Karl Bringmann, Christian Ikenmeyer, and Jeroen Zuiddam. On Algebraic Branching Programs of Small Width. J. ACM, 65(5):1–29, 2018. (Preliminary version in the 32nd Computational Complexity Conference (CCC’17). doi:10.1145/3209663.
\bibitem{ref8} Peter Bürgisser. The complexity of factors of multivariate polynomials. Foundations of Computational Mathematics, 4(4):369–396, 2004. arXiv:1812.06828.
\bibitem{ref9} Peter Bürgisser. On defining integers and proving arithmetic circuit lower bounds. Computa- tional Complexity, 18(1):81–103, 2009. doi:10.1007/s00037-009-0260-x.
\bibitem{ref10} Peter Bürgisser, Michael Clausen, and Amin Shokrollahi. Algebraic complexity theory, volume 315. Springer Science \& Business Media, 2013.
\bibitem{ref11} David V Chudnovsky and Gregory V Chudnovsky. On expansion of algebraic functions in power and Puiseux series, I. Journal of Complexity, 2(4):271–294, 1986. URL: https: //www.sciencedirect.com/science/article/pii/0885064X86900063.
\bibitem{ref12} Michael Coons. The Transcendence of Series Related to Stern’s Diatomic Sequence. Interna- tional Journal of Number Theory, 06, November 2011. doi:10.1142/S1793042110002958.
\bibitem{ref13} Michael Coons and Peter Borwein. Transcendence of power series for some number theoretic functions. Proceedings of the American Mathematical Society, 137, July 2008. doi:10.1090/ S0002-9939-08-09737-2.
\bibitem{ref14} Wellington De Melo and Benar Fux Svaiter. The cost of computing integers. Proceedings- American Mathematical Society, 124:1377–1378, 1996. URL: https://www.ams.org/journals/ proc/1996-124-05/S0002-9939-96-03173-5/S0002-9939-96-03173-5.pdf.
\bibitem{ref15} Richard A. Demillo and Richard J. Lipton. A probabilistic remark on algebraic program testing. Information Processing Letters, 7(4):193–195, 1978. URL: https://www.sciencedirect.com/ science/article/abs/pii/0020019078900674.
\bibitem{ref16} Pranjal Dutta. Real tau-Conjecture for sum-of-squares: A unified approach to lower bound and derandomization. In 16th International Computer Science Symposium in Russia (CSR 2021), 2021. URL: https://drive.google.com/file/d/1X8eo9GM4SCNsC2vWjPbUwMX0vff5i2k3/ view.
\bibitem{ref17} Pranjal Dutta, Nitin Saxena, and Amit Sinhababu. Discovering the roots: Uniform closure results for algebraic classes under factoring. In Proceedings of the 50th Annual ACM SIGACT Symposium on Theory of Computing, pages 1152–1165, 2018. URL: https://www.cse.iitk. ac.in/users/nitin/papers/factor-closure.pdf.
\bibitem{ref18} Pranjal Dutta, Nitin Saxena, and Thomas Thierauf. A Largish Sum-Of-Squares Implies Circuit Hardness and Derandomization. In 12th Innovations in Theoretical Computer Science Conference (ITCS 2021), volume 185 of Leibniz International Proceedings in Informatics (LIPIcs), pages 23:1–23:21. Schloss Dagstuhl–Leibniz-Zentrum für Informatik, 2021. URL: 10.4230/LIPIcs.ITCS.2021.23.
\bibitem{ref19} Joshua A Grochow et al. Complexity in ideals of polynomials: questions on algebraic complexity of circuits and proofs. Bulletin of EATCS, 2(130), 2020. URL: http://bulletin.eatcs.org/ index.php/beatcs/article/view/607.
\bibitem{ref20} Joshua A. Grochow, Ketan D. Mulmuley, and Youming Qiao. Boundaries of VP and VNP. In 43rd International Colloquium on Automata, Languages, and Programming (ICALP 2016), volume 55, pages 34:1–34:14, 2016. URL: https://core.ac.uk/download/pdf/62922137.pdf.
\bibitem{ref21} Gorav Jindal. On approximate polynomial identity testing and real root finding. PhD thesis, Saarland University, 2019. doi:10.22028/D291-29880.
\bibitem{ref22} Erich Kaltofen. Uniform closure properties of p-computable functions. In Proceedings of the eighteenth annual ACM symposium on Theory of computing, pages 330–337, 1986. doi: 10.1145/12130.12163.
\bibitem{ref23} Erich Kaltofen. Single-factor Hensel lifting and its application to the straight-line complexity of certain polynomials. In Proceedings of the 19th annual ACM symposium on Theory of computing (STOC’87), pages 443–452, 1987. doi:10.1145/28395.28443.
\bibitem{ref24} Pascal Koiran. Valiant’s model and the cost of computing integers. computational complexity, 13(3):131–146, 2005.
\bibitem{ref25} Pascal Koiran. Shallow circuits with high-powered inputs. Innovations in Computer Sci- ence (ICS), 2011. URL: https://hal-ens-lyon.archives-ouvertes.fr/ensl-00477023v4/ document.
\bibitem{ref26} Pascal Koiran and Sylvain Perifel. Interpolation in Valiant’s theory. Computational Complexity, 20(1):1–20, 2011. doi:10.1007/s00037-011-0002-8.
\bibitem{ref27} Pascal Koiran, Natacha Portier, Sébastien Tavenas, and Stéphan Thomassé. A $\tau$ -Conjecture for Newton Polygons. Foundations of computational mathematics, 15(1):185–197, 2015. doi: 10.1007/s10208-014-9216-x.
\bibitem{ref28} FV Kuhlmann. On convergent power series, 1996. URL: https://www.mathi.uni-heidelberg. de/~roquette/KONVPOTREIHEN.pdf.
\bibitem{ref29} Hsiang Kung and Joseph Traub. All Algebraic Functions Can Be Computed Fast. J. ACM,
\bibitem{ref30} Dick Lipton and Ken Regan. Factoring and factorials, February 2009. URL: https://rjlipton. wordpress.com/2009/02/23/factoring-and-factorials/.
\bibitem{ref31} Richard J Lipton. Polynomials with 0-1 coefficients that are hard to evaluate. SIAM Journal on Computing, 7(1):61–69, 1978. Preliminary version in the 16th Annual Symposium on Foundations of Computer Science (FOCS 1975). URL: https://epubs.siam.org/doi/abs/ 10.1137/0207004?journalCode=smjcat.
\bibitem{ref32} Richard J Lipton. Straight-line complexity and integer factorization. In International Algorithmic Number Theory Symposium (ANTS 94), pages 71–79. Springer, 1994. doi: 10.1007/3-540-58691-1\_45.
\bibitem{ref33} Richard J Lipton and Larry J Stockmeyer. Evaluation of polynomials with super- preconditioning. Journal of Computer and System Sciences, 16(2):124–139, 1978. URL: https://www.sciencedirect.com/science/article/pii/0022000078900417.
\bibitem{ref34} Meena Mahajan. Algebraic Complexity Classes. In Perspectives in Computational Complexity, pages 51–75. Springer, 2014. doi:10.1007/978-3-319-05446-9\_4.
\bibitem{ref35} Carlos Moreira. On asymptotic estimates for arithmetic cost functions. Proceedings of the American Mathematical Society, 125(2):347–353, 1997. URL: https://www.jstor.org/ stable/2161660.
\bibitem{ref36} Ketan D Mulmuley. The GCT program toward the P vs. NP problem. Communications of the ACM, 55(6):98–107, 2012. doi:10.1145/2184319.2184341.
\bibitem{ref37} Danny Nguyen and Igor Pak. Complexity of short generating functions. In Forum of Mathematics, Sigma, volume 6. Cambridge University Press, 2018. arXiv:1702.08660.
\bibitem{ref38} Øystein Ore. Über höhere kongruenzen. Norsk Mat. Forenings Skrifter, 1(7):15, 1922.
\bibitem{ref39} Igor Pak. Complexity problems in enumerative combinatorics. In Proceedings of the Inter- national Congress of Mathematicians – Rio de Janeiro 2018. Vol. IV. Invited lectures, pages 3153–3180. World Sci. Publ., Hackensack, NJ, 2018. doi:10.1142/9789813272880\_0176.
\bibitem{ref40} David A Plaisted. New NP-hard and NP-complete polynomial and integer divisibility problems. Theoretical Computer Science, 31(1-2):125–138, 1984. Preliminary in the 17th Annual Symposium on Foundations of Computer Science (FOCS 1976). URL: https: //www.sciencedirect.com/science/article/pii/0304397584901300.
\bibitem{ref41} Jacob T Schwartz. Fast probabilistic algorithms for verification of polynomial identities. Journal of the ACM (JACM), 27(4):701–717, 1980. doi:10.1145/322217.322225.
\bibitem{ref42} Adi Shamir. Factoring numbers in O (logn) arithmetic steps. Information Processing Let- ters, 8(1):28–31, 1979. URL: https://www.sciencedirect.com/science/article/abs/pii/ 0020019079900875.
\bibitem{ref43} Amir Shpilka and Amir Yehudayoff. Arithmetic Circuits: A survey of recent results and open questions. Foundations and Trends® in Theoretical Computer Science, 5(3–4):207–388, 2010. doi:10.1561/0400000039.
\bibitem{ref44} Michael Shub and Steve Smale. On the intractability of Hilbert’s Nullstellensatz and an algebraic version of “NP $\neq$ P?”. Duke Math. J., 81(1):47–54 (1996), 1995. A celebration of John F. Nash, Jr. doi:10.1215/S0012-7094-95-08105-8.
\bibitem{ref45} Volker Strassen. Vermeidung von divisionen. Journal für die reine und angewandte Mathematik, 264:184–202, 1973.
\bibitem{ref46} L Valiant. Reducibility by algebraic projections in: Logic and algorithmic. In Symposium in honour of Ernst Specker, pages 365–380, 1982.
\bibitem{ref47} Leslie G Valiant. Completeness classes in algebra. In Proceedings of the 11th Annual ACM symposium on Theory of computing, pages 249–261. ACM, 1979. doi:10.1145/800135.804419.
\bibitem{ref48} Joachim Von Zur Gathen and Jürgen Gerhard. Modern computer algebra. Cambridge university press, 2013.
\bibitem{ref49} Wact. Some Accessible Open Problems. Workshop on Algebraic Complexity Theory (WACT 2016).   URL: https://www.cs.tau.ac.il/~shpilka/wact2016/concreteOpenProblems/ openprobs.pdf.
\bibitem{ref50} Klaus W Wagner. The complexity of combinatorial problems with succinct input representation. Acta informatica, 23(3):325–356, 1986. doi:10.1007/BF00289117.
\bibitem{ref51} Richard Zippel. Probabilistic Algorithms for Sparse Polynomials. In Proceedings of the International Symposium on Symbolic and Algebraic Computation, EUROSAM ’79, pages 216–226, 1979. doi:10.1007/3-540-09519-5\_73.
\bibitem{ref52} J.von zur Gathen and V. Strassen. Some polynomials that are hard to compute. Theoretical Computer Science, 11(3):331–335, 1980. URL: http://www.sciencedirect.com/science/ article/pii/0304397580900201.
\end{thebibliography}
\end{document}
